\chapter{Investigaciones Previas Relacionadas}
\label{chap:investigaciones-previas}

\insertminitoc
\parindent0pt

La revisión de investigaciones previas permite situar el presente trabajo dentro del estado
del arte en reconocimiento de lengua de señas y en el desarrollo tecnológico vinculado al
\acrshort{lesho} en Honduras. Este capítulo organiza los antecedentes en tres áreas: los sistemas de
reconocimiento basados en dispositivos físicos, los sistemas basados en visión por
computadora, y los recursos y desarrollos específicos para el contexto hondureño. A partir
de este análisis se identifican las brechas que fundamentan la propuesta de este proyecto.

\section{Sistemas de Reconocimiento de Lengua de Señas Basados en Sensores Físicos}

Los primeros enfoques para el reconocimiento automático de lengua de señas se
fundamentaron en el uso de dispositivos de captura de datos montados directamente sobre
el cuerpo del usuario. Este tipo de sistemas emplean guantes instrumentados con sensores
inerciales, acelerómetros y giroscopios que miden la posición, orientación y movimiento de
cada dedo y de la mano en su conjunto.

Ji et al. diseñaron un guante de datos de bajo costo equipado con 16 unidades de medición
inercial capaces de capturar movimientos complejos de lengua de señas china. Utilizando
modelos de aprendizaje automático como árboles de decisión, máquinas de soporte
vectorial y bosques aleatorios, junto con una red \acrshort{lstm} con mecanismo de atención,
lograron una precisión de 98.85\,\% en el reconocimiento de 20 señas dinámicas \cite{ji_dataglove_2023}.

Otros enfoques han empleado dispositivos comerciales como el Leap Motion, un sensor
óptico que captura la posición tridimensional de los dedos sin contacto físico. Mittal et al.
desarrollaron un modelo \acrshort{lstm} modificado para el reconocimiento continuo de lengua de
señas utilizando datos del Leap Motion, reportando resultados competitivos en la
clasificación de secuencias gestuales \cite{mittal_lstm_2019}.

Si bien los sistemas basados en sensores alcanzan altas tasas de precisión, presentan
limitaciones prácticas que restringen su adopción fuera del entorno de laboratorio. El costo
de los dispositivos, la necesidad de calibración, la incomodidad de uso prolongado y la
dependencia de hardware especializado constituyen barreras relevantes, especialmente en
contextos de bajos recursos donde la accesibilidad económica es un factor determinante \cite{sadik_survey_2025}.

\section{Sistemas de Reconocimiento Basados en Visión por Computadora}

Como alternativa a los sistemas basados en sensores, el enfoque basado en visión por
computadora utiliza cámaras convencionales para capturar los gestos del usuario y
procesarlos mediante algoritmos de aprendizaje profundo. Este enfoque no requiere
dispositivos adicionales, lo que reduce el costo y mejora la comodidad del usuario.

Los avances recientes en frameworks de estimación de pose, particularmente MediaPipe
de Google, han simplificado la extracción de puntos de referencia anatómicos de la mano
en tiempo real. MediaPipe Hands identifica 21 landmarks por mano con coordenadas
tridimensionales, lo que permite representar cualquier gesto como un vector numérico que
puede ser procesado por modelos de clasificación \cite{zhang_mediapipe_2020}.

La combinación de MediaPipe con redes recurrentes de tipo \acrshort{lstm} se ha consolidado como
uno de los enfoques más utilizados en la literatura reciente para el reconocimiento de señas
dinámicas. Khartheesvar et al. aplicaron MediaPipe Holistic junto con una red \acrshort{lstm} para
el reconocimiento de Lengua de Señas India, alcanzando una precisión de 94.8\,\% en el
dataset INCLUDE y 87.4\,\% en el dataset INCLUDE-50. Los autores destacaron la capacidad
de MediaPipe para capturar secuencias de movimiento dinámico y la eficiencia de \acrshort{lstm}
para modelar las dependencias temporales inherentes a las señas \cite{khartheesvar_indian_2024}.

Srivastava et al. propusieron un sistema de reconocimiento continuo de lengua de señas
que utiliza MediaPipe Holistic para la extracción de landmarks de mano, cuerpo y rostro, y
una red \acrshort{lstm} como clasificador. Evaluado sobre un dataset propio de Lengua de Señas
India con 45 señas y gestos, el sistema alcanzó una precisión de 88.23\,\% en tiempo real \cite{srivastava_continuous_2024}.

La integración de estos modelos en aplicaciones móviles también ha sido explorada.
Sousa et al. desarrollaron una aplicación móvil en Flutter para la traducción del
alfabeto de la Lengua de Señas Portuguesa, integrando dos enfoques entrenados sobre
un dataset propio: una red neuronal convolucional optimizada con TensorFlow Lite para
clasificación de imágenes on-device, y un método basado en MediaPipe para la extracción
de landmarks con clasificación mediante un perceptrón multicapa. La evaluación reportó
valores de F1 de aproximadamente 76\,\% para el enfoque basado en \acrshort{cnn} y 77\,\% para el
enfoque de MediaPipe con \acrshort{mlp}, lo que demuestra la viabilidad de ejecutar el reconocimiento
de señas directamente en el dispositivo y sin conexión a internet \cite{sousa_lgp_flutter_2025}.

A pesar de estos avances, una revisión exhaustiva publicada por Sadik et al. identificó que
la gran mayoría de los sistemas de reconocimiento reportados en la literatura se concentran
en un número reducido de lenguas de señas con grandes conjuntos de datos disponibles,
como la Lengua de Señas Americana (\acrshort{asl}), la Lengua de Señas China (\acrshort{csl}) y la Lengua
de Señas India (\acrshort{isl}). Las lenguas de señas regionales o con menor documentación
permanecen desatendidas en la investigación, lo que limita la transferibilidad de los
resultados a contextos lingüísticos distintos \cite{sadik_survey_2025}.

\section{Antecedentes sobre \acrshort{lesho} en Honduras}

Además de los avances registrados a nivel internacional, resulta necesario examinar el
estado del \acrshort{lesho} desde el punto de vista de su documentación formal y del desarrollo
tecnológico que se ha impulsado en Honduras. Este contexto local permite identificar con
mayor precisión los recursos disponibles y las necesidades específicas que la investigación
busca atender.

El Lenguaje de Señas Hondureño (\acrshort{lesho}) cuenta con un proceso de documentación que se
inició en 1983 con la publicación del primer diccionario impreso denominado ``Las Manos
Hablan'', elaborado en el seno de la Asociación de Sordos de Honduras (\acrshort{ash}). A este le
sucedió una segunda obra en 2005, titulada ``Comuniquemos Mejor'', y un tercer diccionario
en 2019, denominado ``Diccionario de \acrshort{lesho} Tomo I, Primera Edición'', que documenta
más de 900 señas \cite{ash_diccionario_2019}.

En el ámbito digital, el Programa de Servicios a Estudiantes con Necesidades Especiales
(\acrshort{prosene}) de la Universidad Nacional Autónoma de Honduras desarrolló el Diccionario
Digital de \acrshort{lesho}, que constituye el primer recurso de este tipo para la lengua. Este
diccionario contiene más de 700 señas en formato de video y fue creado a partir de la
recopilación de términos del entorno universitario por parte de la comunidad estudiantil
sorda de la \acrshort{unah} \cite{unah_diccionario_2024}.

En cuanto al desarrollo tecnológico, el antecedente más relevante en Honduras es el
Proyecto Sherly, desarrollado por la investigadora Yeny Carías en la \acrshort{unah}. Este sistema
convierte la voz captada por un micrófono en texto, y posteriormente traduce el texto a
una animación de un avatar tridimensional que ejecuta las señas correspondientes en
\acrshort{lesho}. El proyecto fue reconocido por MIT Technology Review y está orientado a reducir
la brecha comunicacional en entornos educativos universitarios \cite{proceso_sherly_2014}.

Williams y Parks realizaron una encuesta sociolingüística de la comunidad sorda hondureña
en la que documentaron la existencia de variaciones regionales del \acrshort{lesho}, el papel central
de la \acrshort{ash} en la estandarización de la lengua, y la situación educativa de las personas sordas
en Honduras. El estudio reportó que más de 3\,000 personas han recibido formación en
\acrshort{lesho} a través de cursos ofrecidos por la \acrshort{ash} y la \acrshort{unah} \cite{williams_encuesta_2015}.

En síntesis, los esfuerzos previos en Honduras se han concentrado en la documentación
lexicográfica del \acrshort{lesho} y en la dirección de traducción de texto o voz hacia lengua de
señas. Durante la revisión bibliográfica realizada no se identificaron antecedentes
documentados de sistemas que aborden la dirección inversa, es decir, el reconocimiento
automático de señas del \acrshort{lesho} mediante visión por computadora para su conversión a
texto o lenguaje natural.

\section{Brechas Identificadas en la Literatura}

A partir de la revisión de los antecedentes presentados en las secciones anteriores, se
identifican las siguientes brechas que justifican el desarrollo de la presente investigación:

En primer lugar, durante la revisión bibliográfica realizada no se identificaron sistemas
de reconocimiento de \acrshort{lesho} basados en visión por computadora. Los trabajos previos en
Honduras se han enfocado exclusivamente en la dirección de texto a señas, dejando sin
atender la dirección inversa que permitiría a una persona sorda expresarse hacia un
oyente sin que este conozca la lengua de señas.

En segundo lugar, durante la revisión bibliográfica realizada no se identificaron
antecedentes documentados de un conjunto de datos de landmarks de \acrshort{lesho} orientado al
entrenamiento de modelos de inteligencia artificial. Los recursos existentes son de
naturaleza lexicográfica, es decir, diccionarios de consulta visual, pero no conjuntos de
datos estructurados para el aprendizaje automático.

En tercer lugar, las soluciones reportadas en la literatura internacional no son
directamente aplicables al contexto hondureño, dado que cada lengua de señas posee un
vocabulario y una gramática propios. Un modelo entrenado con la Lengua de Señas
Americana o la Lengua de Señas India no reconoce señas del \acrshort{lesho}.

En cuarto lugar, los sistemas existentes a nivel internacional operan predominantemente
en una sola dirección. Son escasos los trabajos que implementan una solución bidireccional
que integre tanto el reconocimiento como la representación visual de señas en una misma
aplicación.

Estas brechas definen el espacio de contribución de la presente investigación, que busca
desarrollar un sistema bidireccional de comunicación basado en reconocimiento automático
de \acrshort{lesho} ejecutado on-device en una aplicación móvil, para el cual no se encontraron
antecedentes documentados durante la revisión bibliográfica realizada.
